% Źródła: Wikipedia: Ferdinand Karl Schweikart, Franz Taurinus, Non-Euclidean geometry.
% https://en.wikipedia.org/wiki/Ferdinand_Karl_Schweikart
% https://en.wikipedia.org/wiki/Franz_Taurinus

\section{Schweikart i Taurinus} % lang-pl
\section{Schweikart e Taurinus} % lang-it
\label{sekcja_schweikart_taurinus}

Na początku XIX wieku pytanie o równoległe zaczną zadawać nie tylko matematycy, lecz także prawnicy. % lang-pl
W tej sekcji poznamy wuja i siostrzeńca, którzy zbudują „geometrię astralną'' i wyślą o niej listy do Gaussa. % lang-pl
All'inizio dell'Ottocento la questione delle parallele la porranno non solo i matematici, ma anche i giuristi. % lang-it
In questa sezione conosceremo uno zio e un nipote che costruiranno una «geometria astrale» e ne scriveranno a Gauss. % lang-it

Ferdinand Karl Schweikart urodzi się 28 lutego 1780 roku w Erbach im Odenwald i zostanie profesorem prawa: w Giessen (1809--1812), Charkowie (1812--1816), Marburgu (1816--1821), a od 1821 roku w Królewcu, gdzie umrze 17 sierpnia 1857 roku. % lang-pl
W 1807 roku opublikuje pracę o teorii równoległych, a w 1818 roku, za pośrednictwem swojego ucznia Christiana Ludwiga Gerlinga, przekaże Gaussowi krótką notatkę o \emph{geometrii astralnej}, w której suma kątów trójkąta jest mniejsza od $180^\circ$. % lang-pl
Ferdinand Karl Schweikart nascerà il 28 febbraio 1780 a Erbach im Odenwald e diventerà professore di diritto: a Giessen (1809--1812), Char'kov (1812--1816), Marburgo (1816--1821) e dal 1821 a Königsberg, dove morirà il 17 agosto 1857. % lang-it
Nel 1807 pubblicherà un lavoro sulla teoria delle parallele e nel 1818, tramite il suo allievo Christian Ludwig Gerling, farà giungere a Gauss una breve nota sulla \emph{geometria astrale}, nella quale la somma degli angoli di un triangolo è minore di $180^\circ$. % lang-it
\index[persons]{Schweikart, Ferdinand Karl}%
\index[persons]{Gerling, Christian Ludwig}%

Bratanek Schweikarta, Franz Adolph Taurinus (1794--1874), studiuje prawo w Heidelbergu, Giessen i Getyndze, a potem zamieszka jako prywatny uczony w Kolonii. % lang-pl
W 1824 roku napisze do Gaussa, który zachęci go do dalszych badań, ale poprosi o zachowanie dyskrecji; Taurinus opublikuje więc w 1825 roku \emph{Theorie der Parallellinien}, a w 1826 roku \emph{Geometriae prima elementa}. % lang-pl
Il nipote di Schweikart, Franz Adolph Taurinus (1794--1874), studierà diritto a Heidelberg, Giessen e Gottinga, e poi vivrà come studioso privato a Colonia. % lang-it
Nel 1824 scriverà a Gauss, che lo incoraggerà a proseguire ma chiederà di mantenere il riserbo; Taurinus pubblicherà così nel 1825 la \emph{Theorie der Parallellinien} e nel 1826 la \emph{Geometriae prima elementa}. % lang-it
\index[persons]{Taurinus, Franz Adolph}%

Taurinus podejmie domysł Lamberta (zobacz sekcję \ref{sekcja_lambert}): zbada geometrię na sferze o urojonym promieniu, którą nazwie logarytmiczno-sferyczną, i wyprowadzi w niej wzory trygonometryczne. % lang-pl
Sam jednak pozostanie przy przekonaniu, że geometria euklidesowa jest wyjątkowa. % lang-pl
Rozczarowany słabym odzewem (choć pochlebnie odpowie mu m.in. Georg Ohm), zniszczy większość egzemplarzy \emph{Geometriae prima elementa}. % lang-pl
Taurinus riprenderà la congettura di Lambert (si veda la sezione \ref{sekcja_lambert}): studierà la geometria su una sfera di raggio immaginario, che chiamerà logaritmico-sferica, e vi ricaverà le formule trigonometriche. % lang-it
Rimarrà tuttavia convinto che la geometria euclidea sia speciale. % lang-it
Deluso dalla scarsa accoglienza (anche se tra l'altro Georg Ohm gli risponderà con favore), distruggerà la maggior parte delle copie della \emph{Geometriae prima elementa}. % lang-it
\index[persons]{Ohm, Georg}%

Halsted uzna Schweikarta za pierwszego, który opublikuje świadomy traktat o geometrii nieeuklidesowej; datę jej pierwszego świadomego stworzenia i nazwania umieści między 1812 a 1816 rokiem. % lang-pl
Halsted riterrà Schweikart il primo a pubblicare un trattato consapevole di geometria non euclidea; la data della sua prima creazione e denominazione consapevole la collocherà tra il 1812 e il 1816. % lang-it
\index[persons]{Halsted, George Bruce}%
